\uuid{KMkj}
\titre{Ensembles}
\auteur{bodin ; croizat ; sacre}
\organisation{exo7}
\licence{CC BY-NC-SA 4.0}
\datecreate{2022-03-08}
\niveau{L1}
\module{Algèbre}
\chapitre{Logique, ensemble, raisonnement}
\sousChapitre{Ensemble}
\difficulte{1}
\theme{Ensembles}
\exo7qcmid{hilisit.logique.018}

\contenu{
\begin{question}
\qclassification{100.02}
Quels sont les ensembles ayant au moins $4$ éléments ?
\begin{answers}
	\bad{$\varnothing$}
	\good{$[0,2] \cap [1,3]$}
	\bad{$\{0,3\} \cap \{1,3\}$}
	\good{$\Nn \setminus \{0,1,2,3\}$}
\end{answers}
\begin{explanations}
$\varnothing$ ne contient aucun élément. \\
$[0,2] \cap [1,3] = [1,2]$ contient une infinité d'éléments. Rappel : $[1,2] = \{x\in\Rr \mid 1 \le x \le 2\}$.\\
$\{0,3\} \cap \{1,3\} = \{3\}$ ne contient qu'un seul élément. \\
$\Nn \setminus \{0,1,2,3\} = \{4,5,6,7,\ldots\}$ contient une infinité d'éléments.
\end{explanations}
\end{question}
}
